\chapter{Towards Reliable Systems of Twinned Systems}\label{ch:approach}
This chapter presents an approach to building reliability in SoTS. Building on the gaps identified in the literature review, it formalizes reliability in the context of constituent lifecycle dynamics and defines explicit levels of reliability under disturbance and reconfiguration. The chapter establishes a structured foundation for modeling, assessing, and preserving SoTS operation.

\section{Reliability in SoTS}
Reliability, in classical dependability theory, is defined as the probability that a system performs its intended function without failure for a specified period of time \cite{avizienis2004basic}. In SoS, this notion must be extended to account for operational independence, dynamic reconfiguration, and emergent behavior among constituent systems. Accordingly, reliability is increasingly reframed in terms of disturbances rather than isolated component faults \cite{ferreira2023towards}. 
Disturbances may arise from constituent failures, communication breakdowns, organizational factors, or independent evolution that misaligns services \cite{ferreira2021reliability}.
Because constituents are interdependent, such disturbances propagate across system boundaries and impact mission fulfillment. Reliability in SoS therefore concerns sustaining system-level capabilities despite ongoing disturbance, evolution, and reconfiguration.



% These challenges are amplified in SoTS.
These disturbance dynamics are amplified in SoTS. DTs depend on continuous and timely information exchange between physical and digital constituents. In an evolving SoS environment, observation availability, timing guarantees, and participation cannot be assumed stable. Delayed or incomplete information introduces uncertainty into the digital layer and may compromise coordinated behavior. In SoTS, reliability therefore depends on maintaining acceptable system operation despite fluctuating participation and imperfect information.
% In SoTS, reliability ultimately concerns sustaining system capabilities even as constituent availability and information quality vary over time.

The systematic literature review (\secref{sec:literature-review-key-takeaways}) confirms this gap. Although reliability is frequently cited as a motivation or quality attribute in SoTS studies, it is rarely modeled explicitly or validated at runtime. Most contributions address reliability indirectly through structural redundancy, communication fallback mechanisms, state persistence, or distributed deployment architectures, rather than treating it as an operational property that must be formally assessed and maintained during execution.




\section{Existing Strategies for Reliable SoTS}
\begin{table*}[t]
\small
\centering
\caption{SoS Reliability Strategies with \citet{hanmer2013patterns} Fault-Tolerance Pattern Families}
\label{tab:sos-reliability-techniques}
\renewcommand{\arraystretch}{1.25}
\begin{tabular}{p{3.5cm} p{3.5cm} p{7.5cm}}
\toprule
\textbf{Pattern Family} & \textbf{SoS Reliability Strategy} & \textbf{Techniques Identified \cite{ferreira2021reliability, uday2015designing, abdelmoumen2026comprehensive}} \\
\midrule

\textbf{Architectural Patterns} 
& Structural Redundancy and Substitution 
& Spatial redundancy; functional redundancy; constituent replication; Substitutable Heterogeneous System Architectures (SHSA) \\


\midrule

\textbf{Error Detection} 
& Runtime Monitoring and Verification 
& Fault Detection and Isolation (FDI); runtime verification using Linear Temporal Logic (LTL) and Computation Tree Logic (CTL); anomaly detection \\

\midrule

\textbf{Error Recovery} 
& Task and Goal Reallocation 
& Multi-objective optimization (e.g., NSGA-II) for task redistribution and mission re-planning; dynamic task reassignment; architectural reconfiguration \\


\midrule

\textbf{Error Mitigation} 
& Estimation and Prediction 
& Kalman filtering; Neural Network  prediction (e.g., Long Short-Term Memory (LSTM)); Time-series forecasting\\

\midrule

\textbf{Fault Treatment} 
& Dependency, Probabilistic, and Hazard Modeling 
& Functional Dependency Network Analysis (FDNA); Bayesian Networks; Markov chains and Continuous-Time Markov Models (CTMM); Petri nets; Monte Carlo simulation; 
\newline
Failure Mode and Effects Analysis (FMEA); Failure Mode, Effects, and Criticality Analysis (FMECA); Fault Tree Analysis (FTA); Event Tree Analysis (ETA); Hazard and Operability Study (HAZOP) \\

\bottomrule
\end{tabular}
\end{table*}

Techniques for preserving reliability in SoS align closely with the fault-tolerance pattern families described by \citet{hanmer2013patterns}. As summarized in \tabref{tab:sos-reliability-techniques}, the approaches identified in major SoS reliability surveys \cite{ferreira2021reliability, uday2015designing, abdelmoumen2026comprehensive} span: architectural patterns (structural organization and redundancy to contain faults), error detection (mechanisms to identify abnormal behavior at runtime), error recovery (strategies to restore acceptable operation after disruption), error mitigation (techniques that limit the impact of errors during operation), and fault treatment (analysis and modeling approaches that identify, diagnose, or eliminate underlying causes).

Existing SoTS studies predominantly emphasize architectural patterns and error mitigation mechanisms. Most implementations rely on structural redundancy and component substitution, including fallback communication layers, state persistence, distributed or cloud-based DT deployment, asynchronous messaging, and data caching during disconnection \cite{acharya2023twins, aziz2022empowering, larsen2024towards, lopez2023modeling}. Error detection is occasionally supported through runtime monitoring, event-triggered replanning, or dynamic rescheduling \cite{park2020digital, villalonga2021decision-making}. In contrast, comparatively few studies incorporate explicit fault treatment approaches such as parametric reliability modeling, uncertainty reasoning, or fault tree analysis (FTA) \cite{kutzke2021subsystem, oquendo2019dealing, saraeian2022digital}. 


\section{Ensuring Reliable Participation and Stable Operation in SoTS}
A foundational principle in SoS architecting is the notion of \emph{stable intermediate forms} \cite{maier1998architecting}: partial configurations must remain operational and capable of delivering useful functionality throughout deployment, evolution, and reconfiguration. Subsequent work reinforces this view, arguing that SoS evolution must proceed through viable configurations that preserve value and maintain service continuity during adaptation \cite{rechtin1991systems, rechtin2010art, mekdeci2014pliability, petitdemange2021design, weyns2013challenges}. 

In practice, however, maintaining such stable intermediate forms is difficult because SoS configurations are not static. Reconfiguration is intrinsic to SoS evolution: constituents may join, leave, degrade, suspend participation, or shift roles over time \cite{fang2020improving, manzano2020dynamic, pradhan2016achieving, mohsin2019modeling}. As participation changes, the operational boundary of the SoTS shifts, introducing structural reconfiguration and uncertainty about constituent state and availability. Reliability must therefore be preserved not within a fixed architecture, but across continuously evolving compositions. Without an explicit model of membership transitions, reliability mechanisms respond to isolated disruptions rather than systematically managing how system configurations change over time.

SoS research has increasingly formalized constituent lifecycle states. \citet{axelsson2019refined} distinguish ignorant, prepared, passive, and active states; \citet{hristozov2023modeling} define a state machine governing safe initialization and withdrawal; \citet{petitdemange2021design} structure reconfiguration through coordinated active-passive transitions; and other work models negotiation states \cite{qin2016contract}, standby and replacement roles \cite{du2024reconfiguration}, or lifecycle-driven situational awareness \cite{svenson2020situation}. 

In SoTS, these transitions directly affect the availability and consistency of digital representations. To preserve stable intermediate forms, participation changes must therefore be bounded, coordinated, and, when necessary, compensated. The remainder of this section derives explicit reliability requirements from constituent lifecycle governance and introduces levels of operational reliability that characterize how effectively a SoTS maintains capability under reconfiguration and evolution. 
